\documentclass[aspectratio=169]{beamer}
\usetheme{Madrid}
\usecolortheme{beaver}
\usepackage[utf8]{inputenc}
\usepackage[T1]{fontenc}
\setbeamertemplate{navigation symbols}{}
\setbeamercolor{title}{fg=white,bg=cyan!60!black}
\setbeamercolor{frametitle}{fg=white,bg=cyan!50!black}
\setbeamercolor{block title}{bg=cyan!15,fg=black}
\setbeamercolor{block body}{bg=cyan!5,fg=black}

\title{G6 — Collision \& Physics}
\subtitle{Hitboxes, bouncing, gravity, knockback}
\author{ByteForge --- Grade 7 --- Game Maker}
\date{}

\begin{document}

\begin{frame}
    \titlepage
\end{frame}

\begin{frame}{Today's Plan}
    \begin{itemize}
        \item Hitbox Check
        \item Bounce \& Gravity
        \item Knockback
        \item[$\bigstar$] BOSS --- Breakout Clone
    \end{itemize}
\end{frame}

\begin{frame}{The distance formula — real Pythagoras}
    \Large
    \texttt{dist = sqrt((x2-x1)\textasciicircum{}2 + (y2-y1)\textasciicircum{}2)} — literally Pythagoras' theorem, applied to two circle centres.
\end{frame}


\begin{frame}{Gravity and bouncing — physics as a loop}
    \Large
    Increase velocity by a small amount EVERY frame to simulate gravity — no complicated physics engine needed.
\end{frame}


\begin{frame}{Combining collision detection with movement}
    \Large
    "Detect, then react" — check for a condition (collision), then DO something as a direct result (bounce away).
\end{frame}


\begin{frame}{BOSS — linear search through an array}
    \Large
    An array stores a whole list under one name — checking each item one at a time, in order, is called a linear search.
\end{frame}

\begin{frame}{Real-World Example}
    \Large
    A basketball's arc and bounce follows the exact same idea — gravity constantly pulling it down (velocity increasing each moment) until it hits the floor and reverses.
\end{frame}

\begin{frame}{Watch}
    \Large
    Suggested video topic: \emph{Collision detection and gravity explained — game physics basics}

    \vspace{1em}
    \normalsize
    Search this on YouTube together, or swap in a specific video you've
    already vetted --- do not project an unverified link cold in class.
\end{frame}

\begin{frame}{Let Them Discover}
    \Large
    Drop a ball (real or a video of one bouncing) and watch closely: does each bounce go exactly as high as the last, or does it lose a little height each time? What in the code (Bounce \& Gravity) causes that?
\end{frame}

\begin{frame}{Your Turn}
    \Large
    Open the student map, find \textbf{G6 — Collision \& Physics}, and work through it
    hands-on --- everything you just saw, they now do themselves.
\end{frame}

\end{document}
